\documentclass[12pt,a4paper]{article}

% ---------------------------------------------------
% GRUNDSPRACHE & SCHRIFT
% ---------------------------------------------------
\usepackage[ngerman]{babel}

\usepackage{fontspec}
\setmainfont{Times New Roman}
\usepackage{enumitem}

% ---------------------------------------------------
% SEITENLAYOUT
% ---------------------------------------------------
\usepackage{geometry}
\geometry{
    left=4cm,
    right=2cm,
    top=4cm,
    bottom=2cm,
    headsep=1cm,
    footskip=2cm
}

% ---------------------------------------------------
% FORMATIERUNG
% ---------------------------------------------------
\usepackage{setspace}
\onehalfspacing
\setlength{\parskip}{6pt}
\setlength{\parindent}{0pt}
% Zeilenumbruch-Notreserve gegen Overfull-Boxen bei langen Komposita
\emergencystretch=1em

% ---------------------------------------------------
% FUSSNOTENFORMATIERUNG (einzeilig, 10pt)
% ---------------------------------------------------
\usepackage{etoolbox}
\appto{\footnotesize}{\fontsize{10pt}{12pt}\selectfont}
\renewcommand{\footnotesize}{\fontsize{10pt}{12pt}\selectfont}

\makeatletter
\renewcommand\@makefntext[1]{%
  \noindent\makebox[1.8em][l]{\@makefnmark}%
  \begin{spacing}{1}{\fontsize{10pt}{12pt}\selectfont #1}\end{spacing}}
\makeatother

\usepackage{titlesec}
\titlespacing*{\section}{0pt}{12pt}{6pt}
\titlespacing*{\subsection}{0pt}{12pt}{6pt}
\titlespacing*{\subsubsection}{0pt}{12pt}{6pt}

\titleformat{\section}
  {\normalfont\bfseries\normalsize}{\thesection}{1em}{}

\titleformat{\subsection}
  {\normalfont\bfseries\normalsize}{\thesubsection}{1em}{}

\titleformat{\subsubsection}
  {\normalfont\bfseries\normalsize}{\thesubsubsection}{1em}{}

% Gliederungstiefe: bis Ebene 3 (x.y.z) nummerieren und ins Inhaltsverzeichnis
\setcounter{secnumdepth}{3}
\setcounter{tocdepth}{3}

\pretocmd{\section}{\clearpage}{}{}

% ---------------------------------------------------
% GRAFIKEN, TABELLEN, VERZEICHNISSE
% ---------------------------------------------------
\usepackage{graphicx}
\usepackage{booktabs}
\usepackage{listings}
\usepackage{xcolor}
\definecolor{codegruen}{rgb}{0.0,0.45,0.0}
\definecolor{codegrau}{rgb}{0.45,0.45,0.45}
\definecolor{codeblau}{rgb}{0.10,0.20,0.65}
\definecolor{codehintergrund}{rgb}{0.97,0.97,0.97}
\lstdefinestyle{pythonstil}{
   language=Python,
   backgroundcolor=\color{codehintergrund},
   basicstyle=\ttfamily\footnotesize,
   keywordstyle=\color{codeblau}\bfseries,
   commentstyle=\color{codegruen}\itshape,
   stringstyle=\color{codegrau},
   numbers=left, numberstyle=\tiny\color{codegrau}, numbersep=8pt,
   frame=single, rulecolor=\color{codegrau},
   breaklines=true, showstringspaces=false, captionpos=b,
   literate={ä}{{\"a}}1 {ö}{{\"o}}1 {ü}{{\"u}}1 {ß}{{\ss}}1
            {Ä}{{\"A}}1 {Ö}{{\"O}}1 {Ü}{{\"U}}1
}
\lstset{style=pythonstil}
\renewcommand{\lstlistingname}{Quellcode}   % Beschriftung "Quellcode 5.1"

% ---------------------------------------------------
% VERZEICHNISFORMATIERUNG (tocloft)
% ---------------------------------------------------
\usepackage{tocloft}
\tocloftpagestyle{fancy}

% Überschriften der Verzeichnisse: fett, 12pt, schwarz
\renewcommand{\cfttoctitlefont}{\bfseries\normalsize}
\renewcommand{\cftlottitlefont}{\bfseries\normalsize}
\renewcommand{\cftloftitlefont}{\bfseries\normalsize}

\setlength{\cftafterlottitleskip}{0pt}
\setlength{\cftafterloftitleskip}{0pt}

\renewcommand{\cftsecfont}{\normalfont}
\renewcommand{\cftsecpagefont}{\normalfont}
\renewcommand{\cftsubsecfont}{\normalfont}
\renewcommand{\cftsubsecpagefont}{\normalfont}
\renewcommand{\cftsubsubsecfont}{\normalfont}
\renewcommand{\cftsubsubsecpagefont}{\normalfont}

\renewcommand{\cftsecdotsep}{\cftdotsep}
\renewcommand{\cftsubsecdotsep}{\cftdotsep}
\renewcommand{\cftsubsubsecdotsep}{\cftdotsep}

\setlength{\cfttabindent}{0pt}
\setlength{\cftfigindent}{0pt}

\renewcommand{\cfttabdotsep}{\cftdotsep}
\renewcommand{\cftfigdotsep}{\cftdotsep}
\renewcommand{\cftdotsep}{1}

% Präfix "Tabelle X:" bzw. "Abbildung X:" in den Verzeichnissen
\renewcommand{\cfttabpresnum}{Tabelle }
\renewcommand{\cfttabaftersnum}{:}
\setlength{\cfttabnumwidth}{5em}
\renewcommand{\cfttabfont}{\normalfont}
\renewcommand{\cfttabpagefont}{\normalfont}

\renewcommand{\cftfigpresnum}{Abbildung }
\renewcommand{\cftfigaftersnum}{:}
\setlength{\cftfignumwidth}{6em}
\renewcommand{\cftfigfont}{\normalfont}
\renewcommand{\cftfigpagefont}{\normalfont}

\renewcommand{\cftbeforetabskip}{12pt}
\renewcommand{\cftbeforefigskip}{12pt}

% ---- FORMELVERZEICHNIS -----------------------------
% Eigenes Verzeichnis via tocloft. Nutzung im Text:
%   \begin{equation} ... \end{equation}
%   \formeleintrag{Zielfunktion der Ridge Regression}
\newlistof{formeln}{equ}{Formelverzeichnis}
\renewcommand{\cftequtitlefont}{\bfseries\normalsize}
\setlength{\cftbeforeequtitleskip}{0pt}
\setlength{\cftafterequtitleskip}{0pt}
\renewcommand{\cftformelnpresnum}{Formel }
\renewcommand{\cftformelnaftersnum}{:}
\setlength{\cftformelnnumwidth}{5em}
\setlength{\cftformelnindent}{0pt}
\renewcommand{\cftformelnfont}{\normalfont}
\renewcommand{\cftformelnpagefont}{\normalfont}
\renewcommand{\cftformelndotsep}{\cftdotsep}
\renewcommand{\cftbeforeformelnskip}{12pt}
\newcommand{\formeleintrag}[1]{%
  \addcontentsline{equ}{formeln}{\protect\numberline{\theequation}#1}}

\usepackage[nottoc]{tocbibind}
\usepackage{float}
\usepackage{tabularx}
\usepackage{longtable}
\usepackage[justification=raggedright,singlelinecheck=false]{caption}
\captionsetup{labelfont=bf, textfont=bf}

% ---------------------------------------------------
% LINKS
% ---------------------------------------------------
\usepackage[hidelinks]{hyperref}
\usepackage{csquotes}

% ---------------------------------------------------
% FUSSNOTEN-FORMATIERUNG
% ---------------------------------------------------
\usepackage[hang,flushmargin]{footmisc}
\setlength{\footnotemargin}{1.5em}
\setlength{\footnotesep}{10pt}

% ---------------------------------------------------
% KOPF-/FUSSZEILEN
% ---------------------------------------------------
\usepackage{fancyhdr}
\pagestyle{fancy}
\fancyhf{}
\renewcommand{\headrulewidth}{0pt}
\fancyhead[C]{\thepage}
\setlength{\headheight}{15pt}
\setlength{\headsep}{20pt}

% ---------------------------------------------------
% LITERATURDATENBANK
% ---------------------------------------------------
\begin{filecontents}{literatur.bib}
@article{Bergmeir2012,
  author  = {Bergmeir, Christoph and Ben{\'i}tez, Jos{\'e} Manuel},
  title   = {On the Use of Cross-Validation for Time Series Predictor Evaluation},
  journal = {Information Sciences},
  year    = {2012},
  volume  = {191},
  pages   = {192--213},
  doi     = {10.1016/j.ins.2011.12.028}
}

@article{Breiman2001,
  author  = {Breiman, Leo},
  title   = {Random Forests},
  journal = {Machine Learning},
  year    = {2001},
  volume  = {45},
  number  = {1},
  pages   = {5--32},
  doi     = {10.1023/A:1010933404324}
}

@inproceedings{Bergstra2012,
  author  = {Bergstra, James and Bengio, Yoshua},
  title   = {Random Search for Hyper-Parameter Optimization},
  journal = {Journal of Machine Learning Research},
  year    = {2012},
  volume  = {13},
  pages   = {281--305}
}

@book{Cameron2013,
  author    = {Cameron, A. Colin and Trivedi, Pravin K.},
  title     = {Regression Analysis of Count Data},
  edition   = {2},
  series    = {Econometric Society Monographs},
  number    = {53},
  year      = {2013},
  publisher = {Cambridge University Press},
  address   = {Cambridge},
  doi       = {10.1017/CBO9781139013567}
}

@inproceedings{Chen2016,
  author    = {Chen, Tianqi and Guestrin, Carlos},
  title     = {{XGBoost}: A Scalable Tree Boosting System},
  booktitle = {Proceedings of the 22nd ACM SIGKDD International Conference on Knowledge Discovery and Data Mining},
  year      = {2016},
  pages     = {785--794},
  publisher = {ACM},
  address   = {New York},
  doi       = {10.1145/2939672.2939785}
}

@article{FernandezDelgado2014,
  author  = {Fern{\'a}ndez-Delgado, Manuel and Cernadas, Eva and Barro, Sen{\'e}n and Amorim, Dinani},
  title   = {Do we Need Hundreds of Classifiers to Solve Real World Classification Problems?},
  journal = {Journal of Machine Learning Research},
  year    = {2014},
  volume  = {15},
  pages   = {3133--3181},
  url     = {https://jmlr.org/papers/v15/delgado14a.html}
}

@inproceedings{Grinsztajn2022,
  author    = {Grinsztajn, L{\'e}o and Oyallon, Edouard and Varoquaux, Ga{\"e}l},
  title     = {Why do Tree-Based Models Still Outperform Deep Learning on Typical Tabular Data?},
  booktitle = {Advances in Neural Information Processing Systems (NeurIPS 2022) -- Datasets and Benchmarks Track},
  year      = {2022},
  url       = {https://proceedings.neurips.cc/paper_files/paper/2022/hash/0378c7692da36807bdec87ab043cdadc-Abstract-Datasets_and_Benchmarks.html}
}

@book{Hastie2009,
  author    = {Hastie, Trevor and Tibshirani, Robert and Friedman, Jerome},
  title     = {The Elements of Statistical Learning: Data Mining, Inference, and Prediction},
  edition   = {2},
  year      = {2009},
  publisher = {Springer},
  address   = {New York}
}

@article{Hoerl1970,
  author  = {Hoerl, Arthur E. and Kennard, Robert W.},
  title   = {Ridge Regression: Biased Estimation for Nonorthogonal Problems},
  journal = {Technometrics},
  year    = {1970},
  volume  = {12},
  number  = {1},
  pages   = {55--67},
  doi     = {10.1080/00401706.1970.10488634}
}

@article{Jennings1999,
  author  = {Jennings, Charles R.},
  title   = {Socioeconomic Characteristics and Their Relationship to Fire Incidence: A Review of the Literature},
  journal = {Fire Technology},
  year    = {1999},
  volume  = {35},
  number  = {1},
  pages   = {7--34},
  doi     = {10.1023/A:1015330931387}
}

@article{Jennings2013,
  author  = {Jennings, Charles R.},
  title   = {Social and Economic Characteristics as Determinants of Residential Fire Risk in Urban Neighborhoods: A Review of the Literature},
  journal = {Fire Safety Journal},
  year    = {2013},
  volume  = {62},
  pages   = {13--19},
  doi     = {10.1016/j.firesaf.2013.07.002}
}

@inproceedings{Lundberg2017,
  author    = {Lundberg, Scott M. and Lee, Su-In},
  title     = {A Unified Approach to Interpreting Model Predictions},
  booktitle = {Advances in Neural Information Processing Systems},
  year      = {2017},
  volume    = {30},
  url       = {https://proceedings.neurips.cc/paper_files/paper/2017/hash/8a20a8621978632d76c43dfd28b67767-Abstract.html}
}

@inproceedings{Madaio2016,
  author    = {Madaio, Michael and Chen, Shang-Tse and Haimson, Oliver L. and Zhang, Wenwen and Cheng, Xiang and Hinds-Aldrich, Matthew and Chau, Duen Horng and Dilkina, Bistra},
  title     = {Firebird: Predicting Fire Risk and Prioritizing Fire Inspections in Atlanta},
  booktitle = {Proceedings of the 22nd ACM SIGKDD International Conference on Knowledge Discovery and Data Mining},
  year      = {2016},
  pages     = {185--194},
  publisher = {ACM},
  address   = {New York},
  doi       = {10.1145/2939672.2939682}
}

@book{Molnar2022,
  author    = {Molnar, Christoph},
  title     = {Interpretable Machine Learning: A Guide for Making Black Box Models Explainable},
  edition   = {2},
  year      = {2022},
  publisher = {Independently published}
}

@article{Probst2019,
  author  = {Probst, Philipp and Wright, Marvin N. and Boulesteix, Anne-Laure},
  title   = {Hyperparameters and Tuning Strategies for Random Forest},
  journal = {WIREs Data Mining and Knowledge Discovery},
  year    = {2019},
  volume  = {9},
  number  = {3},
  pages   = {e1301},
  doi     = {10.1002/widm.1301}
}

@article{Rudin2019,
  author  = {Rudin, Cynthia},
  title   = {Stop Explaining Black Box Machine Learning Models for High Stakes Decisions and Use Interpretable Models Instead},
  journal = {Nature Machine Intelligence},
  year    = {2019},
  volume  = {1},
  number  = {5},
  pages   = {206--215},
  doi     = {10.1038/s42256-019-0048-x}
}

@article{Schroeer2021,
  author  = {Schr{\"o}er, Christoph and Kruse, Felix and G{\'o}mez, Jorge Marx},
  title   = {A Systematic Literature Review on Applying CRISP-DM Process Model},
  journal = {Procedia Computer Science},
  year    = {2021},
  volume  = {181},
  pages   = {526--534},
  doi     = {10.1016/j.procs.2021.01.199}
}

@article{Shmueli2010,
  author  = {Shmueli, Galit},
  title   = {To Explain or to Predict?},
  journal = {Statistical Science},
  year    = {2010},
  volume  = {25},
  number  = {3},
  pages   = {289--310},
  doi     = {10.1214/10-STS330}
}

@article{ShwartzZiv2022,
  author  = {Shwartz-Ziv, Ravid and Armon, Amitai},
  title   = {Tabular Data: Deep Learning is Not All You Need},
  journal = {Information Fusion},
  year    = {2022},
  volume  = {81},
  pages   = {84--90},
  doi     = {10.1016/j.inffus.2021.11.011}
}

@inproceedings{Wirth2000,
  author    = {Wirth, R{\"u}diger and Hipp, Jochen},
  title     = {{CRISP-DM}: Towards a Standard Process Model for Data Mining},
  booktitle = {Proceedings of the 4th International Conference on the Practical Applications of Knowledge Discovery and Data Mining},
  year      = {2000},
  pages     = {29--39},
  address   = {Manchester}
}

@incollection{Marban2009,
  author    = {Marb{\'a}n, {\'O}scar and Mariscal, Gonzalo and Segovia, Javier},
  title     = {A Data Mining \& Knowledge Discovery Process Model},
  booktitle = {Data Mining and Knowledge Discovery in Real Life Applications},
  editor    = {Ponce, Julio and Karahoca, Adem},
  publisher = {I-Tech Education and Publishing},
  address   = {Vienna},
  year      = {2009},
  pages     = {1--16},
  doi       = {10.5772/6438}
}
\end{filecontents}

% ---------------------------------------------------
% BIBLATEX-KONFIGURATION (Chicago, einheitlich)
% ---------------------------------------------------
\usepackage[
    backend=biber,
    style=chicago-notes,
    doi=false,
    isbn=false,
    url=false
]{biblatex}
\addbibresource{literatur.bib}

\DefineBibliographyStrings{ngerman}{
    andothers = {et\addabbrvspace al\adddot},
}
\renewcommand*{\finalnamedelim}{\addspace\&\addspace}

\setlength{\bibitemsep}{0.5\baselineskip}
\setlength{\bibhang}{2em}
\renewcommand*{\bibfont}{\normalfont\singlespacing}

\DeclareFieldFormat{postnote}{S.~#1}
\DeclareFieldFormat{multipostnote}{S.~#1}

\renewcommand{\thefootnote}{\arabic{footnote}}

\makeatletter
\renewcommand\@makefntext[1]{%
  \noindent\makebox[1.8em][l]{\@makefnmark}#1%
}

\DeclareNameAlias{sortname}{family-given}
\ExecuteBibliographyOptions{maxnames=2, minnames=1}

\DeclareCiteCommand{\footcite}
  {}
  {%
    \footnote{%
      \printfield{prenote}%
      \iffieldundef{prenote}{}{\addspace}%
      \printnames[labelname]{labelname}\addspace%
      \mkbibparens{\printfield{year}}%
      \iffieldundef{postnote}
        {}
        {\addcomma\space\printfield{postnote}}%
    }%
  }
  {\multicitedelim}
  {}

\makeatother

\AtBeginDocument{%
  \renewcommand{\cfttabfont}{\normalfont}%
  \renewcommand{\cfttabpagefont}{\normalfont}%
  \renewcommand{\cftfigfont}{\normalfont}%
  \renewcommand{\cftfigpagefont}{\normalfont}%
}

% ---------------------------------------------------
% DOKUMENTBEGINN
% ---------------------------------------------------
\begin{document}

% ---------------------------------------------------
% TITELBLATT
% ---------------------------------------------------
\newgeometry{left=2.5cm, right=2.5cm, top=4cm, bottom=2cm}
\begin{titlepage}
\centering

\includegraphics[width=5cm]{FOM_Logo01_2012_rgb.jpg} \\[0.5cm]

\textbf{FOM Hochschule für Oekonomie \& Management} \\
Hochschulzentrum Bonn \\[1.0cm]

Berufsbegleitender Studiengang Wirtschaftsinformatik \\
zum \\
Bachelor of Science (B.Sc.) \\[1.5cm]

Bachelorarbeit zum Thema \\[1.5cm]

{\Large \textbf{„Vorhersage von Feuerwehreinsätzen auf Basis struktureller und
sozioökonomischer Merkmale -- ein Vergleich ausgewählter Machine-Learning-Methoden
am Beispiel von Stadtteildaten“}}

\vfill

\begin{flushleft}
\begin{tabular}{@{}ll}
1. Gutachter: & Prof. Dr. Julian Schröter \\
2. Gutachter: & Oliver Bach \\
Name: & Lukas Nießen \\
Matrikelnummer: & 728085 \\
Wortzahl: & XXXXX Wörter \\
Abgabedatum: & 07.10.2026 \\
\end{tabular}
\end{flushleft}

\end{titlepage}
\restoregeometry

% ---------------------------------------------------
% DISCLAIMER
% ---------------------------------------------------
%\newpage
%\pagenumbering{Roman}
%\setcounter{page}{2}
%\phantomsection
%\section*{Disclaimer}
%\addcontentsline{toc}{section}{Disclaimer}

%\textbf{KI-Nutzung:} In dieser wissenschaftlichen Arbeit wurden zur sprachlichen Unterstützung bei der Erstellung des Textes KI-Modelle eingesetzt. Insbesondere wurden Anthropic Claude und Google Gemini zur Verbesserung der Grammatik, Struktur, Rechtschreibung und Lesbarkeit verwendet. Die inhaltliche Konzeption, Argumentation und jedwede Schlussfolgerungen liegen jedoch vollständig in der Verantwortung des Autors. Sollten Fremdwerke verwendet worden sein, sind diese entsprechend durch Fußnoten gekennzeichnet.

%\textbf{Geschlechterinklusive Sprache:} Darüber hinaus wird in der vorliegenden Arbeit das generische Maskulinum zur besseren Lesbarkeit verwendet. Die in dieser Arbeit verwendeten Personenbezeichnungen beziehen sich, sofern nicht anders kenntlich gemacht, auf alle Geschlechter.

% ---------------------------------------------------
% VERZEICHNISSE
% ---------------------------------------------------
\newpage
\tableofcontents

\newpage
\listoffigures

\newpage
\phantomsection
{\bfseries\normalsize Abkürzungsverzeichnis}
\vspace{6pt}
\addcontentsline{toc}{section}{Abkürzungsverzeichnis}

\begin{tabular}{@{}ll}
ACS & American Community Survey \\[12pt]
AUROC & Area Under the Receiver Operating Characteristic \\[12pt]
CRISP-DM & Cross-Industry Standard Process for Data Mining \\[12pt]
EDA & Explorative Datenanalyse \\[12pt]
MAE & Mean Absolute Error \\[12pt]
RMSE & Root Mean Square Error \\[12pt]
SFFD & San Francisco Fire Department \\[12pt]
SFPD & San Francisco Police Department \\[12pt]
SHAP & SHapley Additive exPlanations \\[12pt]
\end{tabular}


\newpage
\listoftables

\newpage
\phantomsection
\addcontentsline{toc}{section}{Formelverzeichnis}
\listofformeln

% ---------------------------------------------------
% TEXTTEIL
% ---------------------------------------------------
\newpage
\pagenumbering{arabic}
\setcounter{page}{1}

% ===================================================
\section{Einleitung}
% Umfang: ca. 3 Seiten
\subsection{Problemstellung und Praxisrelevanz}
% Einstieg über das Praxisproblem: Ressourcenplanung und Präventionssteuerung
% der Feuerwehr, nicht über die Algorithmen.

% ---- ÜBERNAHME AUS EXPOSÉ (Kap. 1, Abs. 1--2) -- Tempus angepasst ----
% VORHER ERGÄNZEN: 1--2 Absätze zum Praxisproblem (Ressourcenplanung,
% Präventionssteuerung). Das Exposé steigt methodenseitig ein, die Thesis darf das nicht.

Für den Verfahrensvergleich werden drei in der Praxis verbreitete Machine-Learning-Algorithmen
herangezogen, die jeweils auf grundlegend unterschiedlichen methodischen Ansätzen basieren:
Ridge Regression, Random Forest und XGBoost. Welches dieser Verfahren für eine konkrete
Datenkonstellation die beste Prognosegüte liefert, lässt sich theoretisch nicht abschließend
bestimmen, da die Eignung eines Algorithmus maßgeblich von den spezifischen Eigenschaften des
jeweiligen Datensatzes abhängt und daher empirisch unter realen Datenbedingungen untersucht
werden muss.\footcite[3133, 3175]{FernandezDelgado2014}

Als Anwendungsfall dient die Vorhersage von Feuerwehreinsätzen in den Stadtteilen von
San Francisco. Der zugrundeliegende Datensatz integriert öffentlich verfügbare Quellen zu einem
stadtteilbezogenen Datenbestand und umfasst Einsatzdaten des San Francisco Fire Department
(SFFD), sozioökonomische Indikatoren der American Community Survey (ACS), Kriminalitätskennzahlen
sowie bauliche Strukturmerkmale. Die Datenanalyse wird strukturiert nach dem Cross-Industry
Standard Process for Data Mining (CRISP-DM) durchgeführt.
% ---- ENDE ÜBERNAHME ----

\subsection{Zielsetzung und Forschungsfragen}
% Zentrale Forschungsfrage + vier Unterfragen (UF1--UF4).

% ---- ÜBERNAHME AUS EXPOSÉ (Kap. 1, Abs. 3 + Forschungsfragen) -- nahezu 1:1 ----
% ACHTUNG: Dieser Absatz war im Exposé mit Fußnote 2 als KI-unterstützt gekennzeichnet
% (Prompt 1). Der Eintrag muss ins KI-Verzeichnis der Thesis mitwandern.

Die Arbeit verfolgt zwei aufeinander aufbauende Zielsetzungen. Zunächst wird geprüft, ob die
sozioökonomischen, kriminalitätsbezogenen und baulichen Faktoren in statistisch nachweisbarem
Zusammenhang mit der Einsatzlast stehen bzw. einen prädiktiven Beitrag leisten. Aufbauend darauf
werden zwei Zielgrößen modelliert: zum einen die Einsatzhäufigkeit als Anzahl Einsätze pro
Stadtteil und Monat (Regression auf aggregierten Ereignishäufigkeiten), zum anderen die Einsatzart
als kategoriale Zielgröße (Klassifikation). Auf der Einsatzhäufigkeit werden Ridge Regression,
Random Forest und XGBoost trainiert, auf der Einsatzart Random Forest und XGBoost. Die
unterschiedliche Verfahrenszahl folgt der Skalierung der jeweiligen Zielgröße und wird in
Kapitel 6 begründet. Für alle Verfahren werden die Hyperparameter optimiert und die
Prognosegüte anhand etablierter Gütemaße -- Root Mean Square Error (RMSE), Mean Absolute
Error (MAE) und R² für die Regression bzw. Macro-F1 und Area Under the Receiver Operating
Characteristic (AUROC) für die Klassifikation -- im Rahmen einer stadtteilweisen
Kreuzvalidierung gegenübergestellt.\footnote{In
Anlehnung an Anthropic Claude (2026), KI-Verzeichnis Prompt 1}

Daraus ergibt sich folgende zentrale Forschungsfrage:

\begin{quote}
Inwiefern lassen sich die Häufigkeit und die Art von Feuerwehreinsätzen in den Stadtteilen
San Franciscos durch sozioökonomische, kriminalitätsbezogene und bauliche Merkmale statistisch
beschreiben und vorhersagen, und welches der drei Verfahren -- Ridge Regression, Random Forest
oder XGBoost -- erzielt dabei die höchste Prognosegüte?
\end{quote}

Daraus leiten sich folgende Unterfragen ab:

\begin{itemize}[leftmargin=*, itemsep=2pt]
  \item Lassen sich statistisch nachweisbare Zusammenhänge zwischen sozioökonomischen,
        kriminalitätsbezogenen und baulichen Merkmalen und der stadtteilbezogenen
        Einsatzhäufigkeit bzw. Einsatzart feststellen?
  \item Wie unterscheiden sich Ridge Regression, Random Forest und XGBoost hinsichtlich ihrer
        Prognosegüte (RMSE/MAE/R² bzw. F1/AUROC) im Vergleich zu naiven Baselines, bewertet
        mittels zeitreihengerechter Kreuzvalidierung?
  \item Wie verhalten sich die Modelle hinsichtlich Trainings- und Inferenzaufwand?
  \item Welche Implikationen ergeben sich aus den Ergebnissen für die Modellauswahl bei
        vergleichbaren tabellarischen Prognoseaufgaben?
\end{itemize}
% ---- ENDE ÜBERNAHME ----

\subsection{Aufbau der Arbeit}
% Gliederung in Fließtext, inkl. Verweis auf CRISP-DM als Rückgrat von Kap. 3--7.
Bei der vorliegenden Arbeit handelt es sich um eine quantitativ-empirische Untersuchung auf Basis
vorhandener Sekundärdaten. Ihr Aufbau orientiert sich am Cross-Industry Standard Process for Data
Mining (CRISP-DM), einem vom Branchensektor und von der eingesetzten Technologie unabhängigen
Vorgehensmodell für Data-Mining-Projekte.\footcite{Wirth2000} Aufgrund seiner Herstellerneutralität
und seiner breiten Anwendung in der Praxis gilt CRISP-DM bis heute als De-facto-Standard für die
Entwicklung von Data-Mining-Projekten.\footcite[5]{Marban2009}

Das Modell gliedert den Lebenszyklus eines Data-Mining-Projekts in sechs aufeinander aufbauende
Phasen: Business Understanding, Data Understanding, Data Preparation, Modeling, Evaluation und
Deployment. Deren Abfolge ist dabei ausdrücklich nicht starr. Zwischen den Phasen bestehen zwar
typische Abhängigkeiten, welche Phase oder welche konkrete Aufgabe jedoch als nächstes zu
bearbeiten ist, hängt vom Ergebnis der jeweils vorangegangenen Phase ab; ein Rücksprung auf frühere
Phasen ist ausdrücklich vorgesehen. Darüber hinaus ist der Prozess zyklisch angelegt, sodass
gewonnene Erkenntnisse neue Fragestellungen auslösen und den Prozess erneut anstoßen
können.\footcite{Wirth2000}

Kapitel 2 steht außerhalb des Prozessmodells und legt die theoretischen und methodischen
Grundlagen der Arbeit. Es ordnet Feuerwehreinsätze als Prognosegegenstand ein, stellt die
Forschungstradition zu sozioökonomischen und baulichen Einflussfaktoren auf das Brandrisiko dar
und arbeitet den Stand der Forschung zum Modellvergleich auf tabellarischen Daten sowie zu
datengetriebenen Modellen im öffentlichen Sektor auf. Darauf folgt die Darstellung der
untersuchten Verfahren, Ridge Regression, Random Forest und XGBoost, sowie der
Zähldatenmodelle, die als interpretierbare Baseline dienen. Das Kapitel schließt mit der
Herleitung der Forschungslücke und der daraus abgeleiteten Erwartungen, die in Kapitel 8 den
empirischen Befunden gegenübergestellt werden.

Kapitel 3 behandelt die Phase \enquote{Business Understanding}, in der die Projektziele und
Anforderungen aus fachlicher Perspektive in eine Data-Mining-Problemstellung überführt
werden.\footcite[5]{Marban2009} Es beschreibt den Anwendungsfall und die fachliche Zielsetzung,
operationalisiert die beiden Zielgrößen, Einsatzhäufigkeit und Einsatzart, und grenzt den
prognoseorientierten Charakter der Arbeit gegenüber einem kausalen Erklärungsanspruch ab.

Kapitel 4 deckt die Phase \enquote{Data Understanding} ab.\footcite[5]{Marban2009} Das Kapitel
stellt zunächst die Datenquellen und die verwendeten Attribute vor. Es folgt eine explorative
Datenanalyse, welche Verteilungen, räumliche und zeitliche Muster sowie die Klassenverteilung der
Einsatzart untersucht und die daraus resultierenden Qualitätsbefunde dokumentiert.

Kapitel 5 bildet die Phase \enquote{Data Preparation} ab, die sämtliche Aktivitäten umfasst, mit
denen aus den Rohdaten der finale Analysedatensatz konstruiert wird.\footcite[6]{Marban2009} Neben
der Bereinigung und der Behandlung fehlender Werte wird insbesondere die zentrale
Entwurfsentscheidung zur räumlich-zeitlichen Aggregation begründet. Das Kapitel schließt mit dem
Feature-Engineering und der Beschreibung des finalen Analysedatensatzes.

Kapitel 6 konzentriert sich auf die Phase \enquote{Modeling}, in der verschiedene
Modellierungsverfahren ausgewählt, angewendet und ihre Parameter auf optimale Werte kalibriert
werden.\footcite[6]{Marban2009} Es beschreibt den Versuchsaufbau einschließlich der
Validierungsstrategie und der technischen Umsetzung, führt die verwendeten Baselines ein und
dokumentiert die Modellkonfiguration sowie die Hyperparameter-Suchräume. Abschließend werden die
verfahrensspezifischen Konfigurationsentscheidungen und Tuning-Befunde für jeden der drei
Algorithmen dargestellt.

Kapitel 7 ist der Phase \enquote{Evaluation} gewidmet.\footcite[6]{Marban2009} Es stellt die
Prognosegüte für beide Zielgrößen anhand der zuvor festgelegten Gütemaße dar, vergleicht die
Verfahren hinsichtlich ihres Trainings- und Inferenzaufwands und ordnet die Modelle mittels
SHAP-Analyse inhaltlich ein. Das Kapitel schließt mit einem Rückblick auf den durchlaufenen
Analyseprozess.

Die sechste Phase, \enquote{Deployment}, ist nicht Gegenstand dieser Arbeit, da sie einen
konkreten organisatorischen Adressaten voraussetzt, der im Rahmen einer Untersuchung auf
öffentlich verfügbaren Sekundärdaten nicht gegeben ist.\footcite[6]{Marban2009} Die Auslassung
steht im Einklang mit dem dokumentierten Forschungsstand, der die Deployment-Phase als denjenigen
Prozessschritt identifiziert, der in der Anwendung von CRISP-DM am seltensten adressiert
wird.\footcite[531--532]{Schroeer2021}

Kapitel 8 diskutiert die Ergebnisse, bewertet die Algorithmen vergleichend, ordnet den
Erklärungsbeitrag der untersuchten Faktoren ein und reflektiert Limitationen sowie die
Übertragbarkeit der Befunde. Kapitel 9 fasst die zentralen Erkenntnisse zusammen und gibt einen
Ausblick auf weiterführenden Forschungsbedarf.

% ===================================================
\section{Theoretische Grundlagen und Forschungsstand}
% Umfang: ca. 16 Seiten
\subsection{Feuerwehreinsätze als Prognosegegenstand}
% Domäne: Einsatzarten, Einsatzlast, sozioökonomische und bauliche
% Einflussfaktoren (Jennings 1999/2013).

% ---- ÜBERNAHME AUS EXPOSÉ (Kap. 2, Abs. 4, erster Teil) -- 1:1 ----

Eine etablierte Forschungstradition hat gezeigt, dass sich die Häufigkeit von Wohnbränden in
städtischen Quartieren systematisch mit sozioökonomischen und baulichen Strukturmerkmalen
verknüpfen lässt. Jennings zeigt in seinem Review, dass Faktoren wie Armut, Arbeitslosigkeit und
schlechte Wohnverhältnisse konsistent mit erhöhter Brandinzidenz korrelieren, unter anderem weil
in einkommensschwachen Haushalten ältere und wartungsärmere Elektrik sowie Baumaterialien
verbreitet sind. Ergänzend identifiziert er einen ökologischen Erklärungsansatz: Stadtteile mit
hoher Bevölkerungsdichte und baulicher Verdichtung weisen strukturell höhere Brandrisiken auf,
unabhängig vom Verhalten einzelner Bewohner.\footcite[7-9]{Jennings1999}

% AUSBAUEN: Jennings (2013) als Aktualisierung; Einsatzarten und Einsatzlast als Begriffe;
% Abgrenzung Brandinzidenz vs. Einsatzlast (Ihre Zielgröße ist breiter als nur Brände!).
% ---- ENDE ÜBERNAHME ----

\subsection{Datengetriebene Prognose und Algorithmenvergleiche}
% Firebird (Madaio et al. 2016); Benchmark-Literatur (Fernández-Delgado et al.
% 2014, Grinsztajn et al. 2022, Shwartz-Ziv & Armon 2022).

% ---- ÜBERNAHME AUS EXPOSÉ (Kap. 2, Abs. 1--3 + Abs. 4 zweiter Teil) -- 1:1, nur Interpretierbarkeit ausgelagert ----

Die Arbeit ist hinsichtlich der analytischen Verfahren dem Bereich Supervised Machine Learning und
Predictive Analytics zuzuordnen. Sie ordnet sich in die aktuelle Diskussion zum Modellvergleich auf
tabellarischen Daten ein, der für die angewandte Datenwissenschaft von zentraler Bedeutung ist:
Empirische Benchmark-Studien zeigen, dass die Eignung eines Verfahrens stark vom jeweiligen
Datensatz abhängt und kein Algorithmus auf allen Problemklassen gleichermaßen überlegen ist. Die
Eignung eines Verfahrens ergibt sich erst aus der Konfrontation mit konkreten Daten. Damit sind
systematische empirische Algorithmenvergleiche kein Methoden-Beiwerk, sondern wesentlicher
Bestandteil seriöser angewandter Forschung.\footcite[3133, 3175]{FernandezDelgado2014}

Anders als bei Bild- oder Sprachdaten, bei denen Deep-Learning-Verfahren dominieren, gelten bei
strukturierten tabellarischen Daten weiterhin baumbasierte Ensemble-Methoden als Stand der Technik.
Aktuelle Benchmark-Studien zeigen, dass Modelle wie XGBoost und Random Forest auf tabellarischen
Daten mittlerer Größe gegenüber Deep-Learning-Modellen überlegen bleiben, und zwar auch nach
umfangreichem Hyperparameter-Tuning und unabhängig von der Behandlung kategorialer
Variablen.\footcite{Grinsztajn2022} In einer komplementären Untersuchung kommen Shwartz-Ziv und
Armon zu dem Ergebnis, dass XGBoost neuere Deep-Learning-Architekturen auf einer Vielzahl von
Datensätzen übertrifft und dabei einen deutlich geringeren Tuning-Aufwand
erfordert.\footcite[84-90]{ShwartzZiv2022}

Trotz dieser allgemeinen Befunde sind die Ergebnisse einzelner Vergleichsstudien jedoch nicht ohne
weiteres übertragbar. Die Performance eines Algorithmus hängt empirisch nachweislich von
Stichprobengröße, Feature-Anzahl, Signalstärke, Multikollinearitätsstruktur und Interaktionseffekten
ab.\footcite{Grinsztajn2022}

Aktuelle Arbeiten zeigen darüber hinaus, dass datengetriebene Modelle einen substantiellen Beitrag
zur operativen Ressourcenplanung von Feuerwehren leisten können: Das Firebird-Projekt der Stadt
Atlanta berechnet mit Random Forest und SVM Brandrisikoscores für über 5.000 Gebäude und erreicht
True-Positive-Raten von bis zu 71\,\%.\footcite[1-2]{Madaio2016}

% AUSBAUEN: Hsu et al. (2025) -- XGBoost auf Feuerwehr-Einsatzdaten. Steht in der
% Literaturübersicht des Exposés, aber ohne Textbeleg. Fehlt bislang in literatur.bib.
% ---- ENDE ÜBERNAHME ----

\subsection{Untersuchte Verfahren}

\subsubsection{Grundlagen des überwachten Lernens}
% Regression vs. Klassifikation, Bias-Varianz-Zerlegung.

\subsubsection{Ridge Regression}
% Zielfunktion, L2-Regularisierung, Multikollinearität (Hoerl & Kennard 1970).

\subsubsection{Random Forest}
% Bagging, Feature-Subsampling, OOB-Fehler (Breiman 2001).

\subsubsection{XGBoost}
% Gradient Boosting, Objective mit Regularisierung (Chen & Guestrin 2016).

\subsubsection{Zähldatenmodelle als Baseline}
% Poisson / Negativ-Binomial, Überdispersion (Cameron & Trivedi 2013).

\subsection{Modelloptimierung und Modellbewertung}

\subsubsection{Hyperparameter-Tuning, Overfitting und Underfitting}
% Random Search (Bergstra & Bengio 2012), RF-Sensitivität (Probst et al. 2019).

% ---- ÜBERNAHME AUS EXPOSÉ (Kap. 3, Abs. 3, letzter Satz) -- Satzfragment, ausbauen ----

% Die besondere Sensitivität von Random Forest gegenüber Hyperparameter-Einstellungen ist gut
% dokumentiert.\footcite[166]{Probst2019}
% -> Als Beleg brauchbar, trägt aber keinen eigenen Abschnitt. Hier gehört die
%    Funktionsweise von Random Search hin (Bergstra & Bengio 2012) sowie
%    Overfitting/Underfitting -- beides liefert das Exposé nicht.
% ---- ENDE ÜBERNAHME ----

\subsubsection{Gütemaße für Regression und Klassifikation}
% RMSE, MAE, R²; F1, AUROC. Hier gehören die Formeln hin.

\subsubsection{Zeitbewusste Validierung}
% Time-Series-Cross-Validation, Data Leakage (Bergmeir & Benítez 2012).

% ---- ÜBERNAHME AUS EXPOSÉ (Kap. 3, Abs. 4, erster Satz) -- Satzfragment ----

% Die Evaluation erfolgt zeit-bewusst mittels Time-Series-Cross-Validation, die bei zeitlich
% strukturierten Daten Datenleckagen aus der Zukunft vermeidet.\footcite[192-193]{Bergmeir2012}
% -> Der Satz beschreibt Ihr Vorgehen und gehört inhaltlich nach 6.1. Hier ist die
%    allgemeine Darstellung des Verfahrens gefragt (Rolling/Expanding Window).
% ---- ENDE ÜBERNAHME ----

\subsection{Interpretierbarkeit von Modellen}
% SHAP (Lundberg & Lee 2017), Zielkonflikt Güte vs. Interpretierbarkeit
% (Rudin 2019, Molnar 2022), erklärend vs. prognostisch (Shmueli 2010).

% ---- ÜBERNAHME AUS EXPOSÉ (Kap. 2, Abs. 3 zweiter Teil + Kap. 3, Abs. 4) -- 1:1 ----

In der Literatur intensiv diskutiert wird zudem der Zielkonflikt zwischen Prognosegüte und
Interpretierbarkeit: Bei Entscheidungen mit hohem Risiko werden inhärent interpretierbare Modelle
gegenüber nachträglich erklärten Black-Box-Modellen häufig bevorzugt; Post-hoc-Erklärungsmethoden
wie SHapley Additive exPlanations (SHAP) oder Permutation Feature Importance können diesen Konflikt
teilweise entschärfen.\footcite{Rudin2019} Shmueli weist zudem auf die methodisch häufig
vernachlässigte Unterscheidung zwischen erklärender und prognostischer Modellierung hin: Ein Modell
mit hoher Erklärungskraft hat nicht automatisch hohe Prognosegüte und umgekehrt. Algorithmenvergleiche,
die diese Aspekte systematisch dokumentieren, sind daher ein wesentlicher Beitrag zur
evidenzbasierten Modellauswahl in der angewandten Praxis.\footcite{Shmueli2010}

% AUSBAUEN: Funktionsweise der Shapley-Werte (Lundberg & Lee 2017, S. 17), Molnar (2022)
% als Überblick. Das Exposé liefert nur die Einordnung, nicht die Methodik.
% ---- ENDE ÜBERNAHME ----

\subsection{CRISP-DM als Vorgehensmodell}
% Sechs Phasen (Wirth & Hipp 2000); Deployment als dokumentierte Lücke
% (Schröer et al. 2021) und Begründung der Abgrenzung.

% ---- ÜBERNAHME AUS EXPOSÉ (Kap. 3, Abs. 1) -- 1:1, ohne den Satz zur Klassenverteilung (-> 4.2) ----

Methodisch folgt die Arbeit dem CRISP-DM-Schema, das als industrie- und technologieunabhängiges
Prozessmodell für Data-Mining-Projekte gilt und auch zwei Jahrzehnte nach seiner Veröffentlichung
als De-facto-Standard angesehen wird.\footcite[1-2]{Wirth2000} CRISP-DM gliedert ein Projekt in
sechs iterativ durchlaufene Phasen: Business Understanding, Data Understanding, Data Preparation,
Modeling, Evaluation und Deployment. Da der Schwerpunkt dieser Arbeit auf dem Algorithmenvergleich
liegt, erhalten insbesondere die Phasen Modeling und Evaluation überproportionales
Gewicht.\footcite[526-532]{Schroeer2021}

% ERGÄNZEN: Begründung, warum die Deployment-Phase nicht Gegenstand ist. Schröer et al.
% identifizieren genau diese Phase als dokumentierte Lücke -- lässt sich elegant zitieren.
% ---- ENDE ÜBERNAHME ----

\subsection{Forschungslücke und abgeleitete Erwartungen}
% Scharnier zwischen Theorie und Empirie. Ableitung der Erwartungen,
% die in 8.1 und 8.2 gegen die Befunde gestellt werden.

% ---- ÜBERNAHME AUS EXPOSÉ (Kap. 7 Ergebniserwartung) -- Tempus angepasst, Exposé-Rahmung entfernt ----

Aus den dargestellten Befunden lassen sich für den vorliegenden Anwendungsfall zwei Erwartungen
ableiten, die in Kapitel 8 den empirischen Ergebnissen gegenübergestellt werden. Methodisch ist zu
erwarten, dass die baumbasierten Verfahren am stärksten abschneiden.\footcite[9]{Grinsztajn2022}

Inhaltlich lässt sich auf Basis der bestehenden Literatur erwarten, dass Stadtteile mit hoher
Armutsquote, niedrigem Haushaltseinkommen und hohem Altbauanteil eine überdurchschnittliche
Einsatzhäufigkeit und ein abweichendes Einsatzart-Profil aufweisen und diese Merkmale in der
SHAP-Analyse als besonders einflussreiche Prädiktoren hervortreten.\footcite[27-28]{Jennings1999}
Ob kriminalitätsbezogene Merkmale darüber hinaus eigenständige Erklärungskraft besitzen, ist eine
offene Frage, die die Arbeit beantworten soll.

% NICHT übernehmen: Absatz 1 Satz 1 („Ziel der Arbeit ist es aufzuzeigen...") und der
% Schlussabsatz („Als Ergebnis soll eine ... Blaupause ... vorliegen") -- reine Exposé-Rhetorik.
% ---- ENDE ÜBERNAHME ----

% ===================================================
\section{Business Understanding}
% Umfang: ca. 3 Seiten -- CRISP-DM Phase 1
\subsection{Anwendungsfall und fachliche Zielsetzung}
% Wozu dient die Prognose? Warum Stadtteil × Monat als Entscheidungseinheit?

% ---- ÜBERNAHME AUS EXPOSÉ (Kap. 1, Abs. 2) -- Dublette zu 1.1, hier fachlich vertiefen ----
% Der Anwendungsfall ist im Exposé nur in zwei Sätzen umrissen. Für 3.1 fast vollständig
% neu zu schreiben: Wozu dient die Prognose operativ? Welche Entscheidung stützt sie?
% Warum ist Stadtteil × Monat die passende Entscheidungseinheit?
% ---- ENDE ÜBERNAHME ----

\subsection{Zielgrößen, Erfolgskriterien und Abgrenzung}
% Einsatzhäufigkeit (Regression) und Einsatzart (Klassifikation);
% prognoseorientiert, kein kausaler Anspruch.

% ---- ÜBERNAHME AUS EXPOSÉ (Kap. 3, Abs. 4, letzte Sätze) -- 1:1 ----

Die Arbeit verfolgt einen explizit prognoseorientierten Ansatz; kausale Schlussfolgerungen werden
dabei nicht beansprucht.

% ERGÄNZEN: Operationalisierung beider Zielgrößen und Erfolgskriterien.
% Der Bezug zu Shmueli (2010) aus 2.5 sollte hier aufgegriffen werden.
% ---- ENDE ÜBERNAHME ----

% ===================================================
% #########################################################################
% ABGRENZUNG DER KAPITEL 4 BIS 8 -- vor dem Schreiben lesen
% #########################################################################
% Die haeufigste Schwaeche in CRISP-DM-Arbeiten ist, dass dieselbe Sache in
% zwei Kapiteln auftaucht: der Befund im Data Understanding und noch einmal
% als Begruendung in der Data Preparation, die Verfahrensbeschreibung in den
% Grundlagen und noch einmal im Modelling. Schroeter hat "keine Dopplungen,
% Tiefe richtig" ausdruecklich angemerkt (27.07.2026), und das Gutachten zum
% Anwendungsprojekt hat "umfangreich, aber nicht fokussiert" abgestraft (R8).
%
% DIE TRENNUNG IN EINEM SATZ JE KAPITEL:
%
%   Kap. 4  Data Understanding  BEFUND    Was ist da, was faellt auf?
%   Kap. 5  Data Preparation    EINGRIFF  Was wurde damit gemacht, warum?
%   Kap. 6  Modelling           VERFAHREN Was wurde angewandt, warum diese?
%   Kap. 7  Evaluation          ERGEBNIS  Wie gut war es?
%   Kap. 8  Diskussion          DEUTUNG   Was heisst das, was gilt nicht?
%
% DREI TESTFRAGEN, wenn unklar ist, wohin ein Absatz gehoert:
%
%   1. "Gaelte der Satz auch, wenn ich die Daten gar nicht verarbeiten
%       wuerde?"                                    -> ja  = Kapitel 4
%   2. "Veraendert dieser Schritt die Datei auf der Platte?"
%                                                   -> ja  = Kapitel 5
%   3. "Passiert das je Fold, innerhalb der sklearn-Pipeline?"
%                                                   -> ja  = Kapitel 6
%
% GRENZFAELLE, konkret entschieden -- diese Liste ist der eigentliche Nutzen
% des Blocks, weil genau hier sonst doppelt geschrieben wird:
%
%   Klassenverteilung der Einsatzart (79/16/3/1,5 %)          -> 4  Befund
%   Zusammenfassung der NFIRS-Serien zu vier Klassen          -> 5  Eingriff
%   Overdispersion / Dispersionsindex der Zielgroesse         -> 4  Befund
%   Wahl der Negative Binomial als Baseline                   -> 5  Festlegung
%   Baseline-WERTE (R2 0,472 usw.)                            -> 5  Auflage
%                                                                  Schroeter:
%                                                                  Baseline in
%                                                                  die Data
%                                                                  Preparation
%   Golden Gate Park hat 45 Einwohner                         -> 4  Befund
%   Ausschluss der drei Parkgebiete                           -> 5  Eingriff
%   ACS erscheint erst ein Jahr spaeter                       -> 4  Befund
%   Publikationsversatz acs_jahr <= Einsatzjahr - 1           -> 5  Eingriff
%   Stadtteil-Split                                           -> 5  steht als
%                                                                  Spalte IN
%                                                                  der Datei
%   Linearitaetspruefung, Residuen, RESET-Test                -> 6  prueft ein
%                                                                  VERFAHREN,
%                                                                  nicht die
%                                                                  Daten
%   VIF / Multikollinearitaet                                 -> 6  betrifft
%                                                                  Ridge
%   StandardScaler, log(1+y), Label-Encoder                   -> 6  laeuft je
%                                                                  Fold
%   class_weight="balanced"                                   -> 6  Modell-
%                                                                  parameter
%   Hyperparameter-Suchraeume und Tuning-Budget               -> 6
%   Fold-Ergebnisse der drei bzw. zwei Verfahren              -> 7
%   Trainings- und Inferenzzeiten                             -> 7
%   Vergleich der Verfahren untereinander, Rangfolge oder     -> 8
%     deren Verweigerung bei ueberlappender Streuung
%   Effektive Stichprobe von 35 Einheiten, Extrapolation      -> 8  Limitation
%
% ZWEI SONDERFAELLE, die erfahrungsgemaess Aerger machen:
%
%  (a) 4.2 EDA gegen 6.2 Eignungspruefung. Beide schauen Verteilungen an, und
%      derselbe Scatterplot koennte in beiden stehen. Die Regel: In 4.2 steht,
%      WIE DIE DATEN BESCHAFFEN SIND ("die Zielgroesse ist rechtsschief,
%      Dispersionsindex X"). In 6.2 steht, OB EIN VERFAHREN DAZU PASST ("weil
%      sie rechtsschief ist, wird Ridge auf log(1+y) geschaetzt; der
%      RESET-Test verwirft die lineare Spezifikation"). Jede Abbildung
%      erscheint GENAU EINMAL -- im Zweifel dort, wo sie ein Argument traegt,
%      also in 6.2. In 4.2 wird dann darauf verwiesen.
%
%  (b) 2.3 Grundlagen der Verfahren gegen 6. Kapitel 2 erklaert, WIE ein
%      Verfahren funktioniert -- allgemein, lehrbuchhaft, ohne diesen
%      Datensatz. Kapitel 6 sagt, WARUM es hier eingesetzt wird und MIT
%      WELCHEN Einstellungen. Keine Wiederholung der Funktionsweise in 6;
%      wer in 6 erklaeren muss, wie Boosting arbeitet, hat es in 2.3
%      versaeumt.
%
% MERKSATZ: Ein Befund gehoert einmal genannt und danach nur noch zitiert.
% Wer ihn beim zweiten Mal erneut herleitet, schreibt das Kapitel doppelt.
% #########################################################################

\section{Data Understanding}
% Umfang: ca. 7 Seiten -- CRISP-DM Phase 2
% Hier steht NUR der Befund. Jede Konsequenz daraus gehoert nach Kapitel 5,
% jede Verfahrenseignung nach 6.2 -- siehe Abgrenzungsblock oben.
\subsection{Datenquellen und Attributübersicht}
% SFFD-Einsatzdaten, SFPD-Kriminalitätsdaten, Parzellen-/Nutzungsdaten, ACS.
% Übersichtstabelle im Text (Attribut | Quelle | Typ | Ebene | Erhebungsjahr).
% Detailbeschreibung in den Anhang -- Bloat-Gefahr!

% ---- ÜBERNAHME AUS EXPOSÉ (Kap. 5, Abs. 2) -- 1:1; der einzige verwertbare Teil aus Exposé-Kapitel 5 ----

Die empirische Grundlage bildet ein integrierter Datensatz aus mehreren öffentlich zugänglichen
Quellen. Rund 700.000 Einsatzdaten des SFFD, Kriminalitätsdaten des San Francisco Police Department
(SFPD) sowie Parzellen- und Nutzungsdaten wurden über das Open-Data-Portal der Stadt San Francisco
bezogen und auf 41 Stadtteile aggregiert. Sozioökonomische Merkmale auf Census-Tract-Ebene
entstammen den Schätzungen der ACS.

% ERGÄNZEN: Attributtabelle, Erhebungszeiträume, Lizenz-/Abrufdatum der Open-Data-Quellen.
% NICHT übernehmen: der Rechercheweg (Google Scholar, Vorwärts-/Rückwärtssuche) und die
% Auswahlkriterien der Literatur -- reine Exposé-Gattung, entfällt ersatzlos.
% ---- ENDE ÜBERNAHME ----

\subsection{Explorative Datenanalyse und Datenqualität}
% Verteilungen, räumliche und zeitliche Muster, Klassenverteilung der Einsatzart.
% Qualitätsprobleme als Befunde der EDA darstellen, nicht getrennt davon:
% fehlende Werte, Ausreißer, Klassenungleichgewicht, Erhebungsversatz ACS.

% ---- ÜBERNAHME AUS EXPOSÉ (Kap. 3, Abs. 1, letzter Satz) -- Tempus angepasst ----

Im Data Understanding wurde zudem die Klassenverteilung der Einsatzart geprüft; bei ausgeprägten
Ungleichgewichten werden seltene Kategorien zusammengefasst oder die Klassifikation auf zwei
Klassen vereinfacht, etwa Brand vs. Nicht-Brand.

% Hinweis: Der zweite Teilsatz beschreibt eine Aufbereitungsentscheidung und gehört
% streng genommen nach 5.1. Hier nur der Befund, dort die Konsequenz.
% ---- ENDE ÜBERNAHME ----

% ===================================================
\section{Data Preparation}
\label{sec:data-preparation}

% =========================================================================
% GLIEDERUNG NEU (Stand 28.07.2026). Die Pipeline wurde vollstaendig
% umgebaut; der frühere Text dieses Kapitels ist inhaltlich ueberholt und
% wurde entfernt. Alle ZAHLEN fuer dieses Kapitel kommen aus docs/03_STAND.md
% -- dort stehen sie, und nur dort. Nicht aus aelteren Fassungen abschreiben.
%
% ABGRENZUNG (Block vor Kapitel 4): Hier steht der EINGRIFF, nicht der
% Befund. Jeder Befund wurde in Kapitel 4 genannt und wird hier nur noch
% zitiert, nicht erneut hergeleitet. Was je Fold in der Pipeline passiert
% -- Skalierung, log(1+y), Klassengewichte -- gehoert nach Kapitel 6.
% Die Baselines stehen hier und nicht in 6 oder 7: Auflage Schroeter vom
% 27.07.2026, sie gehoeren in die Data Preparation.
% Reihenfolge folgt der Pipeline: prep/s1_daten.py -> s2_datensaetze.py
% -> s3_baselines.py
% =========================================================================

Die Data-Preparation-Phase überführt den in Kapitel~4 beschriebenen
Rohdatenbestand in zwei modellfertige Analysedatensätze. Beide liegen auf
derselben Analyseeinheit -- \emph{Stadtteil~$\times$~Monat} --; der eine misst
die \emph{Menge} der Einsatzlast, der andere ihre \emph{Zusammensetzung}. Diese
Festlegung ist die zentrale Entwurfsentscheidung des Kapitels: Auf
Einzeleinsatz-Ebene würden sich die nur auf Stadtteilebene variierenden
Prädiktoren vielfach wiederholen und Pseudo-Signale erzeugen; eine Aggregation
auf Census-Tract~$\times$~Jahr scheitert daran, dass die Einsatzdaten keine
Tract-Zuordnung tragen, und würde zugleich die Saisonalität eliminieren.

Ein Grundsatz trägt sämtliche Schritte: Jeder Prädiktor muss zum
Prognosezeitpunkt tatsächlich verfügbar gewesen sein.\footcite[192-193]{Bergmeir2012}
Er bestimmt die Auswahl der Spalten ebenso wie die Konstruktion der
Verlaufsmerkmale und die Aufteilung in Trainings- und Testmengen.

% =========================================================================
\subsection{Datenauswahl und Bereinigung}
\label{sec:auswahl-bereinigung}

Dieser Abschnitt beschreibt, welche Daten überhaupt in die Analyse gelangen und
welche bewusst ausgeschlossen werden.

% INHALT:
% - Strikte Spaltenauswahl: von 23 Einsatzspalten bleiben 6. Ausgeschlossen
%   werden Sachschaden, Loeschfahrzeuge, Loeschkraefte, Rettungsdiensteinheiten,
%   Alarmstufe, zivile Tote/Verletzte, Flammenausbreitung -- sie stehen erst
%   NACH dem Einsatz fest und waeren Leakage im engeren Sinn. Argument: Was
%   nicht in den Datensatz kommt, kann nicht versehentlich ins Modell geraten.
% - Bereinigung: 269 doppelte Einsatznummern (0,04 %) entfernt -> 719.989
%   Einsaetze; Ausrueckzeit auf 0-60 min begrenzt (~1,7 % entfallen);
%   Baujahre ausserhalb 1800-2025 als fehlend markiert.
% - Ausschluss ganzer Analyseeinheiten (41 -> 35 Stadtteile):
%     * Golden Gate Park, Lincoln Park, McLaren Park -- Parkgebiete mit 45 bis
%       507 Einwohnern; jede Pro-Kopf-Groesse wird dort beliebig gross
%       (Median der uebrigen Stadtteile: 14.435 Einwohner)
%     * Treasure Island, Lakeshore -- in KEINEM ACS-Jahrgang enthalten
%     * Mission Bay -- erst ab ACS 2021; ein Zutritt mitten in der Zeitreihe
%       erzeugte ein unbalanciertes Panel
%   Zu betonen: Unterschied zwischen zeilenweisem dropna und dem Ausschluss
%   ganzer Einheiten.
% - Analysezeitraum 2015-01 bis 2025-12, hart in config.py fixiert
% ZU FORMULIEREN -- der Start ist eine KONSEQUENZ, keine Setzung:
%   "Der Analysebeginn 2015-01 folgt nicht aus einer Setzung, sondern aus der
%    Publikationsrealitaet des ACS."
%   Herleitung: Die Regel lautet acs_jahr <= Einsatzjahr - 1. Verfuegbar sind
%   die Jahrgaenge 2009, 2014, 2019, 2021, 2023. Ein Einsatz aus 2014 bekaeme
%   damit ACS 2009 - und der enthaelt die Tabelle B15003 nicht, aus der die
%   akademikerquote_pct stammt. Erst ab 2015 greift ACS 2014 (Decision Log #11).
%   NICHT als Wahl darstellen: Es gab keine. und nicht
%   aus den Daten abgeleitet (sonst wandert die Analyse mit jedem Download,
%   und ein angebrochener Randmonat geraet unbemerkt ins Testfenster).
% KEIN Code-Listing.

% =========================================================================
\subsection{Zusammenführung der vier Datenquellen}
\label{sec:integration}

Dieser Abschnitt bleibt bewusst knapp: Er beschreibt, wie die vier Quellen mit
drei unterschiedlichen Raumbezügen und zwei Zeitrastern zu einer Tabelle auf
Einsatzebene zusammengeführt werden.

% INHALT (knapp halten, KEIN Code-Listing):
% - Join A, ACS: zwei Stufen. Census-Tract -> Stadtteil ueber den Crosswalk
%   (Mediane bevoelkerungsgewichtet, Zaehler summiert); dann ZEITBEWUSST mit
%   Publikationsversatz: acs_jahr <= Einsatzjahr - 1. Ohne den Versatz bekaeme
%   ein Einsatz aus 2023 den Jahrgang 2023, der erst Ende 2024 erschien.
%   Als Formel im Fliesstext, nicht als Quellcode.
%   Limitation: Trefferquoten je Jahrgang 2009: 63,1 %; 2014/2019: 79,7 %;
%   2021/2023: 99,2 % (Census-Tract-Grenzen aendern sich, Crosswalk von 2020).
% - Join B, Kriminalitaet: zwei SFPD-Quellen mit Systembruch 05/2018 (CABLE ->
%   Crime Data Warehouse). Die aeltere Quelle hat keine Stadtteilspalte, dafuer
%   Koordinaten -> Spatial Join gegen DIESELBE Neighborhood-Geometrie wie bei
%   Land Use. Sonst bezoegen sich Kriminalitaets- und Baumerkmale desselben
%   Stadtteils auf unterschiedliche Flaechen.
% - Join C, Land Use: 154.544 Parzellen, Zuordnung ueber den Parzellen-
%   Mittelpunkt, Match-Rate 99,5 %. Snapshot 2020, einziger Jahrgang ->
%   zeitkonstantes Strukturmerkmal.
% - Join-Hygiene: Schluesseleindeutigkeit vor jedem Merge, validate-Parameter,
%   Match-Quoten, Zeilenzahl vor/nach gegen den Erwartungswert. Schutz vor
%   kartesischen Produkten.
% ERGEBNIS: einsaetze.parquet, 719.989 x 50, ein Einsatz je Zeile.

% =========================================================================
\subsection{Konstruktion der Merkmale und Zielgrößen}
\label{sec:merkmale}

Dieser Abschnitt ist einer der beiden Schwerpunkte des Kapitels. Er zeigt, wie
aus den zusammengeführten Rohgrößen die Prädiktoren und die vier Zielgrößen
entstehen.

% INHALT (Schwerpunkt, Code-Listings 1 und 2):
% - Quoten aus Zaehlvariablen (Armuts-, Akademiker-, Leerstands-, Altbau-,
%   Wohnnutzungs-, Risikogewerbe-Anteil). Nenner <= 0 ergibt NaN statt
%   Division durch Null.                          -> LISTING 1: safe_ratio
% - Kriminalitaetsindex als Location Quotient, MIT FORMEL:
%     rate(i,t)     = Delikte(i, Fenster endend in t-1) / Einwohner(i)
%     rate(Stadt,t) = Delikte(Stadt, gleiches Fenster) / Einwohner(Stadt)
%     index(i,t)    = rate(i,t) / rate(Stadt,t)
%   Zwei Begruendungen: (a) warum relativ -- der Systembruch 2018 wirkt
%   multiplikativ auf Zaehler und Nenner und kuerzt sich heraus; verbleibende
%   Limitation ist eine Verschiebung der ZUSAMMENSETZUNG der Delikte.
%   (b) warum das Fenster im Vormonat endet -- sonst erklaert Kriminalitaet in
%   t die Einsaetze in t: Beschreibung statt Prognose.
% - Exposure: log_bevoelkerung statt roher Einwohnerzahl. ACHTUNG, die frueher
%   verwendete Begruendung war falsch: Nicht die Armutsquote wechselt das
%   Vorzeichen (+0,49 absolut / +0,46 pro Kopf), sondern die BEVOELKERUNG
%   (+0,20 / -0,42).
% - Aggregation auf Stadtteil x Monat: vollstaendiges Raster, einsatzfreie
%   Monate als echte Nullen; nur ffill, KEIN bfill (Zukunfts-Imputation).
% - Saison als sin/cos statt Monat 1-12 (Dezember und Januar haetten sonst
%   den Abstand 11).
% - Lags: lag_1, lag_12, rolling_mean_3; shift(1) VOR rolling(3); Lag-Vorlauf
%   von 12 Monaten, damit lag_12 schon fuer Januar 2015 definiert ist.
%   WICHTIG: Sie sind KEIN Modellmerkmal mehr. Unter dem Stadtteil-Split
%   (Abschnitt 5.4) waere lag_1 die eigene Vergangenheit des Teststadtteils.
%                                                 -> LISTING 2: Lag-Bildung
% - Zielgroessen (vier):
%     anzahl_einsaetze       Zaehldaten; Mittel 75,9, Median 53, Max 451,
%                            Dispersionsindex 62,8, Nullanteil 0,02 %
%                            -> begruendet NegBin statt Poisson
%     einsaetze_je_1000_ew   Mittel 5,71, Median 3,54, Max 67,7
%     dominante_einsatzart   4 Klassen: Fehlalarm 79,0 %, Techn. Hilfe 16,3 %,
%                            Rettung/EMS 3,1 %, Brand 1,5 %
%     anteil_brand u.a.      stetig in [0,1]; im Mittel 11,6 / 14,9 / 26,2 /
%                            47,4 %
% - Begruendung der Klassifikationsebene (methodisch staerkster Abschnitt):
%   Auf Einzeleinsatz-Ebene tragen alle Einsaetze eines Stadtteil-Monats
%   identische Strukturmerkmale -- 350.481 Zeilen enthielten nur 4.619
%   verschiedene Profile. Ein PERFEKTES Modell haette dort 49,9 % Treffer
%   erreicht gegenueber 48,2 % fuer blosses Raten: Obergrenze 1,7 Prozent-
%   punkte. Auf Stadtteilebene ist die Frage beantwortbar. Die Zielgroesse
%   bleibt eine echte Klasse (argmax ueber die vier NFIRS-Gruppen), keine
%   kuenstliche Einteilung einer stetigen Groesse.\footcite{Altman2006}

% =========================================================================
\subsection{Analyserahmen und Baselines}
\label{sec:rahmen-baselines}

Der zweite Schwerpunkt. Er legt fest, wie die Datensätze in Trainings- und
Testmengen zerfallen, und bestimmt die Referenzwerte, an denen sich jedes
spätere Modell messen lassen muss.

% INHALT (Schwerpunkt, Code-Listing 3):
% - Warum der Split hierher gehoert und nicht nach Kapitel 6: Die Aufteilung
%   steht als Spalten fold und ist_holdout IN DEN DATEIEN. Sie ist eine
%   Eigenschaft des Datensatzes, nicht der Algorithmen -- damit ist die
%   Fairness-Regel konstruktiv abgesichert statt behauptet.
% - STADTTEIL-Split statt Zeitschnitt: 5 Folds a 6 Stadtteile, 6 Stadtteile
%   Hold-out; jeder Stadtteil genau einmal Testfall, mit allen 132 Monaten.
%   Begruendung (zentrale methodische Aussage der Arbeit): Ein Zeitschnitt
%   prueft die Forschungsfrage NICHT, weil dort jeder Stadtteil in Training
%   und Test steht und das Modell sein Niveau bereits kennt.
%   Beleg: 92,5 % der Varianz liegen zwischen den Stadtteilen; der
%   Stadtteil-Mittelwert allein erreicht dort R2 0,888.
% - Doppelte Stratifizierung: sortiert nach brand-dominierten Monaten, bei
%   Gleichstand nach Bevoelkerung. Von 70 brand-dominierten Monaten liegen 35
%   allein in Bayview Hunters Point; ohne Stratifizierung hatte in drei von
%   vier Aufteilungen ein Fold NULL Brand-Testfaelle. Jetzt 13, 9, 6, 3, 2.
%   Kein Leakage -- vgl. StratifiedGroupKFold.  -> LISTING 3: Fold-Zuteilung
% - Modelltauglichkeit: alle Merkmale float64, keine fehlenden Werte. Die
%   Int64-Falle erwaehnen (eine nullable Spalte macht X.to_numpy() zum
%   object-Array; sklearn faengt das still ab, XGBoost lehnt es ab).
% - Baselines (Auflage Schroeter 27.07.2026: gehoeren in die Data Preparation):
%     anzahl_einsaetze       Negative Binomial   R2  0,472 +/- 0,368
%     anzahl_einsaetze       Gesamtmittelwert    R2 -0,832
%     einsaetze_je_1000_ew   Gesamtmittelwert    R2 -2,122
%     anteil_*               Gesamtmittelwert    R2 -0,06 bis -0,11
%     dominante_einsatzart   Mehrheitsklasse     Macro-F1 0,220 / Acc. 0,786
%   Zwei Punkte muessen in den Text:
%     (a) Negative R2 sind korrekt und aussagekraeftig -- wer fuer einen
%         unbekannten Stadtteil den Gesamtdurchschnitt vorhersagt, liegt
%         schlechter als dessen eigener Mittelwert. Genau diese Luecke sollen
%         die Strukturmerkmale schliessen.
%     (b) Die naive Vormonats-Baseline entfaellt: Sie wuerde die eigene
%         Vergangenheit des Teststadtteils nutzen.
%   Zur Auflage "nichtlineare Baseline" (Schroeter 27.07.2026) -- HIER IST
%   EHRLICHKEIT NOETIG: Die frueher gegebene Begruendung (Vormonatswert und
%   saisonaler Durchschnitt seien nichtparametrisch) traegt nicht mehr. Der
%   Vormonatswert ist mit dem Stadtteil-Split entfallen, siehe (b) oben.
%   Gesamtmittelwert und saisonaler Durchschnitt benutzen KEIN EINZIGES
%   MERKMAL und fielen nachgerechnet praktisch zusammen -- faktisch eine
%   triviale Referenz, nicht zwei.
%   Gewaehlt wurde deshalb die NEGATIVE BINOMIAL als alleinige Referenz der
%   Regression (Decision Log #32). Sie ist kein Strohmann (dieselben zwoelf
%   Merkmale, dieselben Folds), ueber den Log-Link auf der Originalskala
%   nichtlinear und verteilungsgerecht fuer Zaehldaten. Entscheidend ist die
%   Trennlinie, die sie zieht: Sie kann Kruemmung, aber KEINE Wechselwirkungen
%   zwischen Merkmalen -- genau die finden RF und XGBoost konstruktionsbedingt.
%   Damit wird der Vergleich zum Beleg statt zur Behauptung.
%   Offen zu benennen: In der Klassifikation gibt es kein Pendant, dort bleibt
%   die Mehrheitsklasse die einzige Referenz und die Latte liegt niedriger.
% - Steckbrief des finalen Datensatzes als Tabelle:
%     Analyseeinheit           Stadtteil x Monat, beide Datensaetze
%     Zeitraum                 2015-01 bis 2025-12 (132 Monate)
%     Beobachtungen Menge      4.620 (35 x 132)
%     Beobachtungen Struktur   4.619 (ein Monat ohne Einsatz)
%     Merkmale                 10 Struktur + 2 Saison = 12
%     Aufteilung               5 Folds a 6 Stadtteile + 6 Hold-out-Stadtteile
%     Ausschluesse             3 Parkgebiete, 3 ohne durchgaengige ACS-Abdeckung
%     Absicherung              19 automatisierte Pruefungen an den Dateien

% ===================================================
\section{Modelling}
% Umfang: ca. 8 Seiten -- CRISP-DM Phase 4
% ABGRENZUNG (Block vor Kapitel 4): Hier steht, WELCHE Verfahren auf den
% fertigen Datensatz angewandt werden und WARUM -- nicht, wie sie
% funktionieren (das steht in 2.3) und nicht, wie gut sie waren (Kapitel 7).
% Alles, was je Fold in der Pipeline laeuft, gehoert hierher: Skalierung,
% log(1+y), Klassengewichte, Label-Encoding, Tuning. Die Eignungspruefung in
% 6.2 prueft VERFAHREN gegen Daten, nicht die Daten selbst -- der Unterschied
% zu 4.2 ist im Abgrenzungsblock unter (a) beschrieben.
\subsection{Versuchsaufbau und technische Umsetzung}
% Der VALIDIERUNGSRAHMEN steht jetzt in Abschnitt 5.4 -- die Aufteilung ist
% eine Eigenschaft des Datensatzes (Spalten fold, ist_holdout) und wird hier
% nur noch referenziert, nicht erneut hergeleitet.
% Hier: identischer Datensatz fuer alle Verfahren, Seeds, wiederholte Splits
% (10 Wiederholungen), Python, Bibliotheken, Versionen. Code in den Anhang.

% ACHTUNG: Der folgende Absatz stammt aus dem Expose und nennt
% Time-Series-Cross-Validation. Das gilt seit dem Umbau NICHT mehr -- der
% Hauptvergleich laeuft ueber einen STADTTEIL-Split (siehe 5.4). Der Absatz
% ist beim Ausformulieren entsprechend zu ersetzen; die Abweichung vom Expose
% gehoert in die Diskussion (8.3) und in die Sprechstunde.

% ---- ÜBERNAHME AUS EXPOSÉ (Kap. 3, Abs. 2 letzter Satz + Abs. 4) -- 1:1, Tempus angepasst ----

Alle Modelle erhalten denselben aufbereiteten Datensatz, um Performance-Unterschiede ausschließlich
auf die Verfahren zurückführen zu können. Die Evaluation erfolgt zeitbewusst mittels
Time-Series-Cross-Validation, die bei zeitlich strukturierten Daten Datenleckagen aus der Zukunft
vermeidet.\footcite[192-193]{Bergmeir2012} Die technische Umsetzung erfolgt in Python.

% ERGÄNZEN: konkrete Fold-Zahl, Fenstergröße, Split-Datum, Seeds, Bibliotheksversionen.
% ---- ENDE ÜBERNAHME ----

\subsection{Eignungsprüfung und Begründung der Verfahrenswahl}
% Die BASELINES sind hierher NICHT mehr zu schreiben -- sie stehen seit dem
% Umbau in Abschnitt 5.4 (Auflage Schroeter, 27.07.2026). Der frueher hier
% stehende Expose-Absatz ist dorthin gewandert.
%
% Dieser Abschnitt beantwortet stattdessen die Frage, warum ueber ein
% einfaches Regressionsmodell hinausgegangen wird. Die Begruendung entsteht
% aus der Regression selbst, in vier Schritten:
%
% 1. Das einfache Modell funktioniert. Ridge auf log(1+y) erreicht R2 0,472
%    und liegt gleichauf mit der Negative-Binomial-Baseline; ein
%    L2-regularisiertes Poisson-GLM erreicht 0,588. Ein lineares Modell ist
%    also ein ernstzunehmender Kandidat, kein Strohmann.
% 2. Der formale Spezifikationstest verwirft es trotzdem. RESET-Test nach
%    Ramsey (1969), H0 "die lineare Spezifikation ist adaequat":
%      Potenzen bis 2: F = 215,23  p = 4e-47   -> H0 verworfen
%      Potenzen bis 3: F = 125,59  p = 4e-53   -> H0 verworfen
% 3. Die Ursache ist lokalisierbar, und sie ist ZWEIFACH (Lauf 03.08.2026):
%    (a) KRUEMMUNG. Von den sechs Merkmalen mit substanzieller Korrelation
%        (|Pearson| > 0,2) zeigt log_bevoelkerung den groessten Abstand
%        zwischen Pearson (+0,416) und Spearman (+0,559) -- ein monotoner,
%        aber gekruemmter Zusammenhang. Wichtig fuer den Text: Bei den
%        uebrigen vier Merkmalen liegt die Korrelation nahe null, dort ist
%        der Pearson-Spearman-Abstand Rauschen und KEIN Befund.
%    (b) WECHSELWIRKUNGEN. Ergaenzt man die 45 Interaktionsterme, steigt das
%        adjustierte R2 von 0,805 auf 0,919.
%    ACHTUNG: Eine fruehere Fassung dieses Kommentars behauptete, die
%    Einzeleffekte seien praktisch linear und nur Wechselwirkungen fehlten.
%    Das ist durch den Lauf vom 03.08.2026 widerlegt und NICHT mehr zu
%    schreiben.
% 4. Daraus folgt die Verfahrenswahl: Ein lineares Modell bildet Kruemmung
%    nur ueber vorab gewaehlte Transformationen ab und Interaktionen nur,
%    wenn man sie von Hand spezifiziert -- bei zehn Merkmalen 45 Terme,
%    willkuerlich ausgewaehlt und bei 23 Trainingsstadtteilen (Fold 1)
%    ueberangepasst. Baumverfahren fangen beides konstruktionsbedingt ab:
%    Ein Split kann an beliebiger Stelle schneiden, und jeder Split bedingt
%    auf die vorherigen.
%
% Ebenfalls hierher: die Linearitaetspruefung als Auflage Schroeter (R7,
% Scatterplots und Residuenanalyse VOR dem Einsatz von Ridge), VIF, sowie die
% Extrapolationspruefung -- unter dem Stadtteil-Split liegen 31 % der
% Testzeilen ausserhalb des Trainings-Wertebereichs, was Baumverfahren und
% Ridge unterschiedlich trifft.
%
% AUFLAGE SCHROETER (E-Mail 03.08.2026): Hierher gehoert ausserdem die
% Begruendung, WARUM Regression und Klassifikation unterschiedlich viele
% Verfahren einsetzen -- drei gegen zwei. Schroeter woertlich: "Sie muessen
% nicht fuer Regression und Klassifikation dieselbe Anzahl oder dieselben
% Algorithmen verwenden. Entscheidend ist, dass die Verfahren zur jeweiligen
% Zielgroesse passen." Die Argumentation, ausformuliert in Decision Log #31:
%
% a) Die Klassifikationsmenge ist eine ECHTE TEILMENGE der Regressionsmenge.
%    Es kommt kein Verfahren hinzu, genau eines faellt weg. Damit entstehen
%    gerade KEINE heterogenen Spezifikationen (Gutachten R1) -- beide
%    Zielgroessen laufen durch denselben Preprocessing- und
%    Validierungsrahmen und werden von denselben Modellfamilien bearbeitet.
% b) RF und XGBoost uebertragen sich, weil bei ihnen nur die Verlustfunktion
%    wechselt (Gini/Entropie statt Varianzreduktion beim Split;
%    multi:softprob statt reg:squarederror). Der Ensemble-Mechanismus ist
%    identisch, beide sind nativ mehrklassenfaehig.
% c) Ridge uebertraegt sich nicht: Es minimiert den quadratischen Fehler auf
%    einer metrischen Zielgroesse. dominante_einsatzart ist nominal mit vier
%    ungeordneten Klassen -- ohne Ordnung und Abstand ist der quadratische
%    Fehler nicht definiert. RidgeClassifier wurde geprueft und verworfen
%    (+-1-Kodierung, One-vs-Rest, keine kalibrierten Wahrscheinlichkeiten).
%    Das Pruefen und Verwerfen gehoert in den Text -- es zeigt, dass die
%    Option gesehen und nicht uebersehen wurde.
% d) Die logistische Regression waere das korrekte lineare Gegenstueck und
%    schlaegt im Vortest die Baseline (Macro-F1 0,282 gegen 0,220 der
%    Mehrheitsklasse, RF 0,301). Sie entfaellt aus FOKUSGRUENDEN, nicht
%    mangels Eignung: ein vierter Verfahrensstrang bei zwei Zielgroessen ist
%    genau die Verbreiterung, die das Gutachten unter R4 als groesstes
%    Notenrisiko benennt. DAS IST SO ZU SCHREIBEN -- eine funktionierende
%    Option als untauglich darzustellen waere unbelegt (R9).

\subsection{Modellkonfiguration und Hyperparameter-Suchräume}
% Tabelle: Verfahren | Hyperparameter | Suchraum | Iterationen | bester Wert.

% ---- ÜBERNAHME AUS EXPOSÉ (Kap. 3, Abs. 3) -- 1:1, Tempus angepasst ----

In der Modeling-Phase wurden Ridge Regression, Random Forest und XGBoost auf beiden Zielgrößen,
Einsatzhäufigkeit (Regression) und Einsatzart (Klassifikation), trainiert. Für jedes Modell erfolgt
ein Hyperparameter-Tuning mittels Randomized Search. Die besondere Sensitivität von Random Forest
gegenüber Hyperparameter-Einstellungen ist dabei gut dokumentiert.\footcite[166]{Probst2019}

% ERGÄNZEN: die Konfigurationstabelle selbst -- das Exposé liefert nur den Rahmen.
% ---- ENDE ÜBERNAHME ----

\subsection{Modellspezifische Tuning-Befunde}
% Auffälligkeiten pro Verfahren -- KEINE Wiederholung der Funktionsweise (steht in 2.3).

% ===================================================
\section{Evaluation}
% Umfang: ca. 11 Seiten -- CRISP-DM Phase 5
%
% =========================================================================
% ERGEBNISINVENTAR -- was am Ende vorliegt und was davon in die Arbeit kommt
% =========================================================================
% Gerechnet werden 400 Einzellaeufe. Sie kommen NICHT in die Arbeit, sondern
% ins Abgabe-Zip; berichtet werden Aggregate.
%
%   Rohdaten in results/          Regression   Klassifikation
%     Einzellaeufe (*_folds.csv)         300              100
%       = Zielgroessen x Verfahren x 10 Wiederholungen x 5 Folds
%     Aggregate (*_mittel.csv)             6                2
%     Tuning-Ergebnisse                   30               10
%     Hold-out                             6                2
%     Baselines je Fold (liegen vor)      20               10
%
%   Berichtete Aussagen          Regression   Klassifikation
%     PRIMAER: gegen Stufe-2-Baseline      6                2
%     sekundaer: paarweise Verfahren       6                1
%     Laufzeitpaare (Training/Inferenz)    6                2
%     Hold-out-Messungen                   6                2
%
% ACHT PRIMAERAUSSAGEN tragen das Kapitel. Jede lautet: "Verfahren X schlaegt
% auf Zielgroesse Y die Stufe-2-Baseline um Z, bei Streuung S ueber zehn
% Wiederholungen." Das ist die Antwort auf Unterfrage 2 (Decision Log #34).
% Jede steht fuer sich - sie haengt nicht davon ab, ob sich die Verfahren
% UNTEREINANDER trennen lassen.
%
% Die sieben paarweisen Vergleiche sind die Rangfolge. Sie kommen nur ins
% Kapitel, wenn der gepaarte Wilcoxon-Test sie hergibt.
%
% EHRLICHKEIT BEI DER TESTZAHL: Insgesamt sind es 15 Signifikanztests - acht
% primaere und sieben paarweise. Die Gesamtzahl gehoert in den Text, sonst
% entsteht der Eindruck, sie sei verschwiegen worden (docs/06_RISIKEN.md, R-10).
%
% MASSGEBLICHE STREUUNG ist std_wiederholungen (ueber die 10
% Wiederholungsmittel), NICHT std_folds ueber alle 50 Laeufe - die 50 sind
% nicht unabhaengig, es sind dieselben 29 Stadtteile in anderer Gruppierung
% (R-5). Beide Spalten stehen in *_mittel.csv.
%
% VIER TABELLEN, DREI ABBILDUNGEN decken alles ab:
%   T1  Aggregate Regression (6 Zeilen)
%   T2  Aggregate Klassifikation (2 Zeilen)
%   T3  Vergleichstabelle mit Wilcoxon-Ergebnissen
%   T4  Hold-out (eine einzige Messung an 6 Stadtteilen, KEIN Mittelwert)
%   A1  Boxplot je Verfahren ueber die 50 Laeufe, je Strang
%   A2  Balken: Verfahren gegen Stufe-2-Baseline -- die Primaeraussage
%   A3  Streudiagramm Laufzeit gegen Prognoseguete -- beantwortet UF3 und UF4
%
% Alle Abbildungen erzeugt modelle/m05_abbildungen.py aus den CSV-Dateien,
% nicht von Hand. Sie sind damit reproduzierbar und bei neuen Laeufen in einem
% Befehl aktualisiert.
% =========================================================================
\subsection{Ergebnisse Einsatzhäufigkeit (Regression)}
% LEITFRAGE: Haengen Einsatzhaeufigkeit und Stadtteilmerkmale zusammen --
% koennen die Modelle die ZAHL der Einsaetze vorhersagen?
%
% Zielgroessen: anzahl_einsaetze (Zaehldaten) und einsaetze_je_1000_ew (Rate).
% Guetemasse RMSE / MAE / R2, je Zeile (Stadtteil x Monat), Mittelwert und
% Streuung ueber alle Folds und Wiederholungen.
% Referenz: die Baselines aus Abschnitt 5.4 -- Negative Binomial R2 0,472;
% Gesamtmittelwert -0,832 (Anzahl) bzw. -2,122 (Rate).
% Zu berichten: Ridge auf log(1+y), Random Forest, XGBoost. Wo sich die
% Streuungsbereiche ueberlappen, ist das so zu schreiben (Gutachten R6).
%                                                          -> UF2

\subsection{Ergebnisse Einsatzart (Klassifikation)}
% LEITFRAGE: Haengen Einsatzart und Stadtteilmerkmale zusammen -- koennen die
% Modelle die ART der Einsaetze vorhersagen?
%
% Zielgroesse: dominante_einsatzart -- vier Klassen je Stadtteil x Monat
% (Brand, Rettung/EMS, Technische Hilfe/Gefahr, Fehlalarm/Good Intent).
% Eine echte Klasse, gebildet als argmax ueber die vier NFIRS-Gruppen des
% Monats; kein gesetzter Schwellwert und keine kuenstliche Einteilung einer
% stetigen Groesse.
%
% Guetemasse: Macro-F1 (Hauptmass) und Macro-AUROC gegen die
% Mehrheitsklassen-Baseline aus 5.4 (Macro-F1 0,220 / Accuracy 0,786).
% ACCURACY IST HIER KEIN HAUPTMASS -- 79 % der Stadtteil-Monate werden von
% Fehlalarm dominiert, ein Modell das immer "Fehlalarm" sagt erreicht 0,786.
%
% Verfahren: NUR Random Forest und XGBoost -- zwei statt drei. Ridge hat auf
% einer nominalen Zielgroesse keine Entsprechung, die logistische Regression
% wurde aus Fokusgruenden gestrichen. Freigabe Schroeter 03.08.2026, begruendet
% in 6.2, festgehalten als Decision Log #31. Im Text ist zu benennen, dass der
% Klassifikationsstrang damit bewusst schmaler ist als der Regressionsstrang
% -- nicht stillschweigend zwei Verfahren berichten, wo das Expose drei
% verspricht.
%
% Klassenverteilung, gehoert in die Ergebnisdarstellung:
%   Fehlalarm 79,0 % | Techn. Hilfe 16,3 % | Rettung/EMS 3,1 % | Brand 1,5 %
% Die duenne Besetzung von Brand ist als Limitation zu benennen: Je Fold
% liegen 2 bis 13 brand-dominierte Monate im Test.
%                                                          -> UF2

\subsection{Trainings- und Inferenzaufwand}
% Laufzeiten, Tuning-Aufwand.                             -> UF3

\subsection{Interpretierbarkeit: SHAP-Analyse}
% Einflussstärke der Faktorgruppen.                       -> UF1

% ---- ÜBERNAHME AUS EXPOSÉ (Kap. 3, Abs. 4) -- Satzfragment ----

% „Zur inhaltlichen Einordnung werden die Modelle mittels SHAP-Werten analysiert, einem
% fundierten Erklärungsansatz für komplexe Modelle."\footcite[17]{Lundberg2017}
% -> Ankündigung, kein Ergebnis. Der Beleg ist verwertbar, der Abschnitt selbst
%    entsteht vollständig neu aus Ihren Analysen.
% ---- ENDE ÜBERNAHME ----

\subsection{Evaluation des Analyseprozesses}
% Rückblick auf den CRISP-DM-Durchlauf; Abgrenzung der Deployment-Phase.

% ===================================================
\section{Diskussion}
% Umfang: ca. 5 Seiten
\subsection{Vergleichende Bewertung der Algorithmen}
% Erwartung (2.7) gegen Befund.
% Zwei Punkte gehoeren zwingend hierher:
% - Spezifikationsasymmetrie: Die NegBin-Baseline erhaelt log(Bevoelkerung)
%   als OFFSET und modelliert damit direkt Einsaetze je Einwohner; Ridge, RF
%   und XGBoost bekommen dieselbe Groesse nur als gewoehnliches Merkmal.
%   Entweder wird der Vorteil ausgeglichen oder offen benannt -- sonst
%   vergleicht die Arbeit nicht drei Algorithmen, sondern drei
%   Spezifikationen.
% - Unterscheidbarkeit: Bei 35 Analyseeinheiten streuen die Fold-Ergebnisse
%   stark. Wo sich die Streuungsbereiche ueberlappen, ist das so zu
%   berichten und nicht als Rangfolge zu kaschieren (Gutachten R6).

\subsection{Erklärungsbeitrag der untersuchten Faktoren}
% Sozioökonomisch / baulich / kriminalitätsbezogen.       -> UF1

\subsection{Limitationen und Übertragbarkeit}
% Methodenkritik, datensatzspezifische Grenzen der Ergebnisse und daraus
% abgeleitete Implikationen für die Modellauswahl.       -> UF4
% ZU FORMULIEREN -- wo das Erkenntnislimit wirklich sitzt:
%   "Eine Ausweitung des Zeitraums wuerde die Aussagekraft nicht erhoehen -
%    die Beschraenkung liegt bei 35 Analyseeinheiten, nicht bei der Zahl der
%    Beobachtungen."
%   Begruendung: Ein zusaetzliches Jahr braechte 35 x 12 = 420 Zeilen (+9 %),
%   aber NULL zusaetzliche Stadtteile. Unter einem Stadtteil-Split besteht die
%   Testmenge aus Stadtteilen; ob ein unbekannter Stadtteil mit 132 oder 144
%   Monaten beschrieben wird, aendert nichts - zumal die Strukturmerkmale
%   innerhalb eines Jahres nahezu konstant sind.
%   Das nimmt der naheliegenden Pruefungsfrage "warum nicht mehr Daten?" die
%   Spitze und zeigt, dass das Limit verstanden ist.
% Aufzunehmen: effektive Stichprobe von 35 Querschnittseinheiten;
% Extrapolation bei unbekannten Stadtteilen (31 % der Testzeilen);
% Hold-out aus nur 6 Stadtteilen; jaehrliche ACS-Merkmale und Land-Use-
% Snapshot 2020; ACS-Trefferquote 2009 von 63,1 %; sowie die Abweichungen
% vom Expose (Stadtteil-Split statt Zeitreihen-CV, Klassifikation mit nur
% zwei statt drei Verfahren, Einsatzart auf Stadtteilebene).
% Zur Verfahrenszahl gehoert hierher der ehrliche Preis der Entscheidung
% (#31): In der Klassifikation fehlt ein lineares Referenzverfahren, die
% Symmetrie linear/Bagging/Boosting besteht also nur im Regressionsstrang.
% Der Vergleich dort ist damit einer zwischen zwei Baumverfahren, die sich
% konstruktionsbedingt aehnlicher sind als Ridge, RF und XGBoost einander.

% ===================================================
\section{Fazit und Ausblick}
% Umfang: ca. 2 Seiten

% ---------------------------------------------------
% LITERATURVERZEICHNIS
% ---------------------------------------------------
\newpage
\nocite{*}  % TODO: entfernen, sobald echte \footcite-Aufrufe im Text stehen
\printbibliography[heading=bibintoc,title={Literaturverzeichnis}]

% ---------------------------------------------------
% ANHANG
% ---------------------------------------------------
\newpage
\phantomsection
\section*{Anhang}
\addcontentsline{toc}{section}{Anhang}

% A1  Vollständige Attributtabelle
% A2  Code-Auszüge (Datenaufbereitung, Tuning, Evaluation)
% A3  Ergänzende Abbildungen (EDA, SHAP)

% ---------------------------------------------------
% KI-VERZEICHNIS
% ---------------------------------------------------
\newpage
\phantomsection
\section*{KI-Verzeichnis}
\addcontentsline{toc}{section}{KI-Verzeichnis}

\begin{table}[H]
    \centering
    \caption*{Verzeichnis der verwendeten KI-Hilfsmittel}
    \vspace{4pt}
    \small
    \begin{tabularx}{\textwidth}{@{}p{0.6cm} p{2.8cm} X@{}}
        \toprule
        \textbf{Nr.} & \textbf{KI-System} & \textbf{Einsatzzweck / Prompt} \\
        \midrule
        \addlinespace
        1 & Anthropic Claude & \enquote{...} \\
        \addlinespace
        2 & Google Gemini & \enquote{...} \\
        \bottomrule
    \end{tabularx}
\end{table}

\end{document}